\section{Availability}

Four public repositories. The first two are the predecessors these notes refer to; the third is the project the examples come from; the fourth is the template \S7 describes.

\begin{itemize}
\item \textbf{\texttt{visgraf/bioeye}} --- a thin vertical slice through the framework, written first. It closed the accumulation loop and not the perception--action loop.
\item \textbf{\texttt{visgraf/active-stereo}} --- a full implementation with verified geometry and a decision record. The loop never ran. Both the failure that motivated the low-ceremony lane and the constitution clause that talked the reviewer out of verifying are in this repository's history.
\item \textbf{\texttt{visgraf/bio-3d-vision}} --- the project. Thirteen experiments, each with its pre-registration, runner, raw results, verdicts and findings; two ledgers; a technical report with its claims-verification pass; and a forward-only history in which every correction appears as an amendment.
\item \textbf{\texttt{visgraf/math-ai-method}} --- the template, instantiable.
\end{itemize}

Every claim these notes make about what happened can be checked in the third. That is the point of having kept it.
